% Methods Section — v2 (2026-05-14)
% Aligned with results_v4.tex:
%   - LORO metric: ForwardEncoding acc_exact (SRM-aligned)
%   - LOCO metric: adjacent accuracy (adj_acc, ±1 hue step) + vulnerability vector
%   - Statistics: Crawford & Howell (1998) for individual CVD; 1,000 per-subject label permutations for above-chance tests; N=300 resamples for model selection
%   - $\Delta$RDM: LOO-consistent, Crawford & Howell for disparity elevation
%   - Sections reordered to match Results narrative (LORO → LOCO → $\Delta$RDM → 2-comp → filter)
% 2026-05-14: Loss replaced (HYBRID → Option C, neural-primary + heavy Tikhonov).
%   Brettel sign-prior claim removed per 2026-05-14 revalidation (Brettel predicts
%   collapse, β_c ≈ 0, not subtype-specific rotation in our cosine parameterization).
%   HYBRID backup at methods_v2_HYBRID_backup.tex.

\section{Methods}
\label{sec:methods}

The analysis follows three stages. First, we characterized each CVD participant's cortical color representation with three neural analyses. Leave-one-run-out (LORO) decoding tested whether color classification is preserved (§\ref{sec:methods:loro}). Leave-one-color-out (LOCO) decoding tested whether continuous hue interpolation is impaired, producing a per-hue vulnerability profile (§\ref{sec:methods:loco}). Representational-geometry analysis measured how far the CVD hue geometry departs from the HC reference, both overall (Procrustes disparity) and pair by pair ($\Delta$RDM) (§\ref{sec:methods:rdm}). Second, two candidate CVD distortion model classes were fit by minimizing a combination of neural and behavioral losses, with parameters selected by a pre-specified three-gate procedure (§\ref{sec:methods:candidates}--§\ref{sec:methods:selection}). Third, the fitted distortion was numerically inverted to yield a per-subject stimulus-space correction filter (§\ref{sec:methods:filter}). Figure~\ref{fig:pipeline} summarizes the fitting and inversion procedure. This filter was then evaluated in a second session of fMRI and behavioral testing, using the hue-discrimination JND and 8AFC color-identification tasks described in §\ref{sec:methods:behavioral}.

% ── Figure 1 ─────────────────────────────────────────────────────────────────
\begin{figure*}[tb]
  \centering
  \includegraphics[width=\textwidth]{fig1_generated_v2}
  \caption{\textbf{Stimuli, task, and analysis pipeline.} \textbf{(A)} Stimulus set: 8 isoluminant CIE $L^{*}a^{*}b^{*}$ hues at 45\degree\ angular spacing ($L^{*} = 75$, chroma $= 40$; red--magenta), shown in the $a^{*}b^{*}$ plane.
  \textbf{(B)} Task: rapid serial visual presentation (RSVP). A fixation period (randomized 3--6\,s) is followed by a 1.5\,s uniform colored disc; 5--9 trials per color per run (Neurodesign-optimized allocation, identical across participants), six runs.
  \textbf{(C)} Analysis pipeline.
  Stage A (Alignment): two-stage GLM single-trial $\beta$ amplitudes, Procrustes-aligned across runs, projected into an SRM space trained on HC participants only.
  Stage B (Neural feature): LORO classification, LOCO interpolation, and the SRM-space representational dissimilarity matrix (RDM).
  Stage C (Modeling \& filter): a 2-component cortical model with S-cone-axis ($\beta_s$) and confusion-axis ($\beta_c$) components; its pre-image yields the per-subject stimulus-space filter.
  HC: $n = 7$ (sub-01--07).}
  \label{fig:paradigm}
\end{figure*}

% ──────────────────────────────────────────────────────────────────────────────
\subsection{Participants}
\label{sec:methods:participants}
% ──────────────────────────────────────────────────────────────────────────────

Thirteen volunteers were recruited under Institutional Review Board approval (\#2510/002-023) and provided written informed consent. Color vision was screened using the 14-plate Ishihara pseudoisochromatic test \cite{ishihara1917}. Participants identifying $\geq 9$ plates correctly were classified as healthy controls (HC; $n = 7$, 3 female, age $22.7 \pm 2.5$ years). Two male participants scoring below this threshold were classified as CVD. The deutan participant was deuteranomalous (age 24 years, 5 of 14 plates correct) and the protan participant was protanomalous (age 25 years, 7 of 14 plates correct). Four additional participants were excluded. Three did not complete the first scanning session, and a fourth completed it but became unavailable for the filter-evaluation session. All CVD results are interpreted as single-case demonstrations using the \citeA{crawford1998} modified $t$-test; no population-level inference is made.

% ──────────────────────────────────────────────────────────────────────────────
\subsection{fMRI stimuli and task}
\label{sec:methods:stimuli}
% ──────────────────────────────────────────────────────────────────────────────

Eight isoluminant chromatic stimuli were drawn from CIE $L^{*}a^{*}b^{*}$ space at equal angular spacing ($0^\circ$--$315^\circ$, $45^\circ$ steps; $L^{*} = 75$, chroma $= 40$) \cite{cie1986}, spanning red, orange, yellow, green, cyan, blue, purple, and magenta (Figure~\ref{fig:paradigm}A). A gray filler ($a^{*} = b^{*} = 0$) served as a null condition.

Stimuli were presented in a Rapid Serial Visual Presentation (RSVP) design to sustain attention and fixation. Each stimulus was a uniform circular chromatic patch (a colored disc, 1.5 s per stimulus) shown in pseudo-random order optimized with Neurodesign \cite{durnez2018}. Participants pressed a button upon detecting two consecutive 'K's within a rapid serial letter stream presented at the center of the chromatic disc. Each run comprised 72 events. Because the schedule was optimized for detection power rather than for equal counts, trials were not balanced across colors within a run: each chromatic stimulus occurred 5--9 times per run (38--47 times across the six runs), with the remaining 14--20 events per run assigned to the gray filler. The same optimized schedule was used for every participant, so this allocation contributes a common component to all individual estimates; single-trial amplitudes were subsequently averaged within run and color, so all downstream analyses operated on a balanced 6 (run) $\times$ 8 (color) design. Inter-stimulus intervals were drawn from $\{3, 4.5, 6\}$ s. Each session comprised six runs ($\approx 7$ min each). The experiment was coded in PsychoPy 2022.2.5 \cite{peirce2019}.

% ──────────────────────────────────────────────────────────────────────────────
\subsection{Behavioral tasks}
\label{sec:methods:behavioral}
% ──────────────────────────────────────────────────────────────────────────────

Two behavioral tasks --- a just-noticeable-difference (JND) hue-discrimination task and an eight-alternative forced-choice (8AFC) color-identification task --- were administered on the same day as the main scanning session (Session~1) for all participants. Both CVD participants additionally repeated both tasks under each filter in the second session (see the \nameref{sec:methods:filter_eval} section).

\paragraph{Just-noticeable difference (JND).}
On each trial participants viewed one hue pair and judged whether the two patches were the same or different. The two hues were separated by a fraction $t$ of the pair's full separation, interpolated in CIE $LCh$, and $t$ served as the staircase level. Two staircases per pair were interleaved from starting levels $t = 0.8$ and $t = 0.5$, both following a 1-up/1-down rule that converges on the $50\%$ ``different'' point \cite{levitt1971}. The step size was $0.15$ for the first two reversals and $0.03$ thereafter. A staircase terminated at eight reversals and at least 20 trials, extended by two further reversals whenever the post-warmup reversal levels had an SD above $0.10$. The threshold was the mean of the last six post-warmup reversal levels.

Eight hue pairs were tested. Three were hypothesized as universally difficult for both CVD types (yellow--purple, blue--purple, red--orange), two probed the deutan confusion axis (orange--yellow, yellow--green), two probed the protan confusion axis (cyan--magenta, green--blue), and one served as a control (red--cyan). For each pair, the JND ratio $\gamma = \gamma_{\rm CVD} / \bar{\gamma}_{\rm HC}$ was computed, where $\bar{\gamma}_{\rm HC}$ is the HC mean threshold. Ratios $\gamma > 1$ indicate elevated hue-discrimination difficulty relative to HC. These ratios enter the fitting loss as described in \S\ref{sec:methods:twocomp}.
% TERMINOLOGY (resolved 2026-08-05): "discrimination" is now reserved for the
% BEHAVIORAL JND construct throughout the paper; the neural LORO construct is
% "color classification". "hue pair" is the unified term (was "color pair" in
% \S\ref{sec:methods:twocomp}).

\paragraph{8AFC color identification.}
On each trial a circular chromatic patch was presented, and participants selected the matching hue from eight options identical to the scanner stimulus set. Identification accuracy served as an independent behavioral readout and did not enter the fitting loss. It was compared to the Ishihara screening outcome post-hoc.

% ──────────────────────────────────────────────────────────────────────────────
\subsection{MRI acquisition and preprocessing}
\label{sec:methods:mri}
% ──────────────────────────────────────────────────────────────────────────────

Data were acquired on a Siemens 3T MAGNETOM Cima.X scanner with BioMatrix 3T coils. To constrain coverage to posterior occipital cortex we used 24 oblique slices perpendicular to the calcarine sulcus \cite{brouwer2009} (TR 1.5 s, TE 30 ms, FA $75^\circ$, $2 \times 2 \times 2$ mm, $96 \times 80$ matrix, 288--292 volumes per run).

Data were converted to BIDS format and defaced with ezBIDS \cite{levitas2024}. Preprocessing followed a custom pipeline rather than a standard container; no confound regression was applied at any stage (Supplementary~\S S1). The T1w image was skull-stripped (FSL \texttt{bet2}; \citeNP{smith2002}), and functional images were registered to it via mutual-information coregistration initialized from the scanner header (FreeSurfer \texttt{mri\_coreg}; \citeNP{wells1996, maes1997}). T1w-to-MNI normalization used 12-DOF affine (FSL FLIRT; \citeNP{jenkinson2002}) followed by nonlinear warping (FNIRT; \citeNP{andersson2007}) onto the MNI152NLin2009cAsym template, yielding 2 mm isotropic MNI-space BOLD data. Each functional volume was resampled once by this composed transform, without slice-timing correction, motion realignment, or spatial smoothing. Head motion was estimated separately, by rigid-body registration of every volume to its run's middle volume (FSL MCFLIRT; \citeNP{jenkinson2002}), and retained as a quality-control measure: mean framewise displacement was $0.32 \pm 0.04$ mm across the nine analyzed participants (HC $0.31$, CVD $0.34$; Supplementary~\S S1). A sensitivity analysis repeating the primary neural endpoints on an independently preprocessed version of the same data is reported in Supplementary~\S S17.
% TODO: Supplementary S17 (preprocessing sensitivity) must exist before submission.
% Scope agreed 2026-08-05: re-run C010 on the true fMRIPrep derivatives (HMC + STC
% + fieldmap SDC applied) and recompute (a) hV4 LOCO adjacent accuracy, (b) SRM
% Procrustes disparity. Primary pipeline stays the method3 custom chain.
% Description reconciled 2026-08-05 against the canonical driver
% analysis/phase0_preprocessing/scripts/run_method3_header_mi_all_subjects.sbatch
% (exp1) and run_method3_header_mi_2nd.sbatch (exp2): bet2 L190-192,
% mri_coreg --regheader L253-266, FLIRT/FNIRT L298-311, template L311.
% Method 2 (header->BBR) and the aCompCor confound path were never adopted;
% their code is in analysis/phase0_preprocessing/_archive/.

% ──────────────────────────────────────────────────────────────────────────────
\subsection{ROI definition and response estimation}
\label{sec:methods:roi}
% ──────────────────────────────────────────────────────────────────────────────

Bilateral V1, V2, V3, and hV4 were defined from the Wang probabilistic atlas \cite{wang2015} at $> 50\%$ probability, intersected with each participant's BOLD brain mask (voxel counts: V1 $665 \pm 209$; V2 $458 \pm 113$; V3 $105 \pm 19$; hV4 $64 \pm 18$; mean $\pm$ SD, $n = 9$). No further functional voxel selection was applied.
% Voxel counts updated 2026-08-03 from the stored C010 amplitudes actually used
% in every analysis: shape[2] of full_dataset_C010/sub-{01..09}/{ROI}/
% amplitudes_procrustes.npy (n = 9; sub-10 excluded from all analyses).
% Superseded values were V1 655+-214; V2 451+-145; V3 103+-29; hV4 63+-22 —
% source and subject set unidentified; SDs differed most (V2 145 vs 113,
% V3 29 vs 19). See docs/PAPER/TODO_additional_analysis.md A9.
% Per-subject counts (sub-01..09):
%   V1  568 405 858 858 858 858 330 560 692
%   V2  402 335 557 557 557 557 258 400 498
%   V3  106  94 115 115 115 115  59 114 115
%   hV4  67  69  70  70  70  70  16  70  70   <- sub-07 hV4 = 16 (known outlier)

Single-trial response amplitudes were estimated with a two-stage GLM \cite{dale1999, brouwer2009, brouwer2013}. First, an ROI-level hemodynamic response function (HRF) was extracted from a finite-impulse-response (FIR) model (8-TR basis set).
% R^2 voxel-selection sentence removed 2026-08-05: it belonged to the superseded
% Baseline32 pipeline, whose orchestration code was archived the same day
% (analysis/comprehensive/_archive/baseline32_orchestration/). C010 applies no
% voxel selection. Full evidence: docs/PAPER/TODO_additional_analysis.md A9.
Second, the estimated HRF and its temporal derivative were convolved with per-color design matrices and fitted by ordinary least squares, alongside per-run linear and constant drift regressors. This yielded one amplitude per (voxel, color, run), organized as a $6 \times 8 \times V$ array per subject.

Voxel amplitudes were aligned across runs using within-subject Procrustes rotation \cite{gower1975}: runs 2--6 were mapped to run 1 via the orthogonal transformation minimizing $\|X_\text{ref} - XQ\|_F$ (scaling prohibited).

% ──────────────────────────────────────────────────────────────────────────────
\subsection{Shared Response Model}
\label{sec:methods:srm}
% ──────────────────────────────────────────────────────────────────────────────

Cross-subject alignment and dimensionality reduction were performed with the Shared Response Model (SRM; \citeNP{chen2015}). SRM estimates per-participant orthonormal mappings $W^{\mathrm{srm}}_i \in \mathbb{R}^{V \times k}$ ($V$ voxels, $k$ latent dimensions) and a group-shared response matrix $S \in \mathbb{R}^{k \times 8}$:

\begin{equation}
  \min_{W^{\mathrm{srm}}_i,\,S}\;\sum_{i=1}^N \|X_i - W^{\mathrm{srm}}_i S\|_F^2 \quad\text{s.t.}\quad {W^{\mathrm{srm}}_i}^\top W^{\mathrm{srm}}_i = I_k
\end{equation}

SRM was trained exclusively on the seven HC participants to avoid circular inference in CVD comparisons. Each CVD participant was subsequently projected into the fixed HC-derived space by solving $W^{\mathrm{srm}}_\text{CVD} = UV^\top$ where $U\Sigma V^\top = \text{SVD}(X_\text{CVD}\cdot S^\dagger)$. Reduced dimensions were chosen by mean-rank aggregation across seven-fold leave-one-subject-out (LOSO) cross-validation, in which one HC is withheld from SRM training on each fold, using three SRM-quality metrics (RDM split-half reliability, cross-subject RDM correlation, and reconstruction error): V1 $k = 4$; V2 $k = 4$; V3 $k = 3$; hV4 $k = 3$.

% ──────────────────────────────────────────────────────────────────────────────
\subsection{Forward encoding model}
\label{sec:methods:encoding}
% ──────────────────────────────────────────────────────────────────────────────

Both classification (LORO) and interpolation (LOCO) decoding used the same forward encoding model \cite{brouwer2009}. The model represents each stimulus as a $F = 6$-dimensional channel vector $\mathbf{c}(\theta) = [c_1(\theta), \ldots, c_F(\theta)]^\top$, where $c_k(\theta) = \max\!\bigl(0,\,\cos(\theta - \theta_k)\bigr)^2$ is a half-wave rectified squared cosine tuned to preferred hue $\theta_k$ (preferred hues spaced equally at $60^\circ$; \citeNP{brouwer2009}). Stacking the eight stimulus channel vectors row-wise gives the channel design matrix $C \in \mathbb{R}^{8 \times F}$. Decoding and encoding optimize different targets and therefore use different estimators of the channel-to-voxel weight matrix $W \in \mathbb{R}^{F \times V}$ (distinct from the SRM mapping $W^{\mathrm{srm}}$).

\paragraph{Decoding (stimulus hue from voxels).}
Decoding recovers the stimulus hue from voxel responses. We estimated $W$ by ordinary least squares, taking the pseudoinverse of the channel design matrix $C$. For a held-out voxel pattern $\mathbf{x}$, channel activations were reconstructed as $\hat{\mathbf{c}} = W\mathbf{x}$. The decoded hue was the angle, among all 360 integer hues, whose basis vector correlated most strongly with $\hat{\mathbf{c}}$ (\citeNP{brouwer2009}; Figure~\ref{fig:forward}). For eight-class exact accuracy this angle was assigned to its nearest stimulus hue. Alternative decoders evaluated during model selection are described in Appendix~A.

\paragraph{Encoding (voxel responses from channels).}
Encoding predicts the voxel responses from the channels ($\hat{X} = CW$, with $\hat{X} \in \mathbb{R}^{8 \times V}$). We estimated $W$ by ridge regression, with the penalty $\alpha$ selected per voxel by generalized cross-validation (GCV; \citeNP{golub1979}; see Supplementary~\S S6). Regularization is warranted because the few stimuli and the correlated basis channels would otherwise overfit $W$ and degrade prediction of held-out responses (\citeNP{kay2008,naselaris2011}). This ridge prior is isotropic and imposes no spatial structure on the voxels.

% ── Figure 2 ─────────────────────────────────────────────────────────────────
\begin{figure*}[!htbp]
  \centering
  \includegraphics[width=\textwidth]{fig_forward_encoder}
  \caption{\textbf{Forward encoding model schematic.}
  Stage 1: each displayed hue is encoded as channel activations by six half-wave rectified squared-cosine basis functions, with preferred hues spaced at 60 degrees (Brouwer \& Heeger 2009).
  Stage 2 (training): the channel-to-voxel weights $W$ are fit by ordinary least squares for hue decoding, or by GCV-regularized ridge regression for voxel-response prediction (encoding).
  Stage 3 (testing): for a held-out voxel pattern, channel activations are reconstructed as $\hat{\mathbf{c}} = W\mathbf{x}$.
  Stage 4: the decoded hue is recovered as the angle, among all 360 integer hues, whose basis vector has the highest Pearson correlation with the reconstructed channel pattern.
  The same model serves both decoding schemes (Section~\ref{sec:methods:loro}): cross-run classification (LORO) and cross-color interpolation (LOCO). Their basis and correlation-based selection are identical, and only the held-out test set differs. Decoding uses the OLS pseudoinverse weights.}
  \label{fig:forward}
\end{figure*}

% ──────────────────────────────────────────────────────────────────────────────
\subsection{Two decoding schemes: cross-run classification and cross-color interpolation}
\label{sec:methods:loro}
% ──────────────────────────────────────────────────────────────────────────────

Two complementary cross-validation schemes used the forward encoding model to characterize each participant's hue processing, differing only in what was held out. Leave-one-run-out (LORO) held out repeated presentations to test the categorical representation of the eight hues, measured as color classification. Leave-one-color-out (LOCO) held out entire hues to test the continuous representation, measured as hue interpolation.

\paragraph{Cross-run classification (LORO)}
Color classification was assessed with leave-one-run-out cross-validation in SRM-aligned space. For each of the six runs, $W$ was estimated on the remaining five runs using mean-subtracted SRM projections, and each held-out run's 8 color patterns were decoded by the correlation readout of \S\ref{sec:methods:encoding}. Primary classification performance was the 8-class exact accuracy (chance $= 1/8 = 0.125$), averaged across six folds.

Group-level HC--CVD comparison used a Mann--Whitney $U$ test on cross-subject generalization scores. The encoding model trained on each HC participant's SRM-aligned data was applied to every other participant's aligned patterns, yielding 21 unique HC-to-HC and 14 HC-to-CVD pairwise scores. Because each participant contributes to several pairs, these scores are not independent, so the test is reported as a descriptive group-level check. Within-ROI single-case tests used the \citeA{crawford1998} modified $t$-test (one-tailed, lower).

\paragraph{Cross-color interpolation (LOCO)}
\label{sec:methods:loco}
Continuous hue interpolation was assessed with leave-one-color-out cross-validation. Each of the 8 hues was held out in turn; $W$ was estimated on the remaining 7 hues using all 6 runs (pooled $42$-sample design matrix), and the held-out hue was predicted by the same correlation readout.

\subparagraph{Adjacent accuracy}
The primary cross-color decoding metric was adjacent accuracy: the proportion of predictions that fell within $\pm 1$ hue step ($45^\circ$) of the target hue, averaged across the 6 runs of each test color (range $0$--$1$; chance level under a uniform random predictor $= 3/8 = 0.375$; exact-accuracy chance $= 1/8 = 0.125$). For HC above-chance testing and for pairwise HC--CVD comparisons at hV4, $p$-values were obtained from 1,000 random color-label permutations and from the \citeA{crawford1998} modified $t$-test (one-tailed, lower), respectively.

\subparagraph{Vulnerability profile}
The 8-dimensional vector of per-hue adjacent-accuracy values, $\mathbf{v} \in [0,1]^8$, constitutes each participant's \emph{hue vulnerability profile}. For group-level model fit, we report Spearman $\rho$ between the observed $\mathbf{v}$ and the $\mathbf{v}$ predicted by the CVD distortion model (see the \nameref{sec:methods:candidates} section). Per-hue group comparisons (Figure~\ref{fig:loco}C) used the \citeA{crawford1998} modified $t$-test (one-tailed, uncorrected). Effect sizes are reported as Hedges' $d$ \cite{hedges1985} with the Crawford--Howell correction for small-$n$ comparisons.

% ──────────────────────────────────────────────────────────────────────────────
\subsection{Representational geometry}
\label{sec:methods:rdm}
% ──────────────────────────────────────────────────────────────────────────────

To characterize the representational geometry in SRM-aligned space, we computed two measures. Procrustes disparity captures how far the CVD geometry lies from the HC reference overall. The per-pair RDM difference ($\Delta$RDM) captures which color pairs drive that departure and is the readout that the inverse fit targets (§\ref{sec:methods:twocomp}).

\paragraph{Procrustes disparity}
Disparity was the Procrustes distance between a participant's SRM-aligned $8 \times k$ pattern and the HC reference. Both patterns were column-centered and scaled to unit Frobenius norm, and disparity was the residual Frobenius norm after the optimal orthogonal rotation. To provide an unbiased HC reference for the held-out HC participant, the HC disparity distribution was estimated by a leave-one-HC-out (LOO) procedure, in which each HC participant is compared to the mean of the remaining six. Under the directional prediction that CVD geometry lies farther from the HC reference than HC geometry does, individual CVD disparity was compared to the HC LOO distribution by the \citeA{crawford1998} modified $t$-test (one-tailed, upper; $\alpha = 0.05$). In this primary estimate the shared SRM space is trained on all seven HC, so each HC reference is in-sample whereas the projected CVD participant is out-of-sample. To control for this asymmetry we additionally computed a symmetric leave-one-subject-out (LOSO) estimate in which the shared space is retrained on six HC and the held-out HC is projected by SVD identically to each CVD participant, with the same Crawford--Howell test applied. The all-HC estimate is primary because every other analysis in the pipeline (LORO, LOCO, $\Delta$RDM, and the inverse fit) operates in that same space; the symmetric LOSO is reported as a disclosed sensitivity analysis (\cref{tab:disparity_loso}).

\paragraph{Per-pair RDM difference ($\Delta$RDM)}
For each participant and ROI, the RDM was the $8 \times 8$ matrix of pairwise correlation distances ($1 - \text{Pearson}\,r$) among the SRM-aligned color patterns. The per-pair difference between each CVD participant's RDM and the HC mean was $\Delta\text{RDM} = \text{RDM}_\text{CVD} - \overline{\text{RDM}}_\text{HC}$, where $\overline{\text{RDM}}_\text{HC}$ is the grand mean RDM over the seven HC participants (the per-pair pattern is shown as a heatmap, Figure~\ref{fig:geometry}A). This per-pair difference is the target of the representational loss $L_{\rm RDM}$ used in the inverse fit (§\ref{sec:methods:twocomp}).

% ──────────────────────────────────────────────────────────────────────────────
\subsection{Candidate CVD distortion models}
\label{sec:methods:candidates}
% ──────────────────────────────────────────────────────────────────────────────

Two candidate model classes were evaluated. Each parameterizes the angular distortion $\delta\theta(\theta)$ that maps the displayed hue angle $\theta$ to the angle perceived by a CVD observer.

% ── Pipeline schematic — inverse-problem workflow ────────────────────────────
\begin{figure*}[tb]
  \centering
  \includegraphics[width=\textwidth]{fig3_workflow}
  \caption{\textbf{Fitting a per-subject cortical distortion and inverting it into a stimulus-space filter.}
  \textbf{(1)} A two-parameter forward model expresses the per-hue distortion as $\delta\theta(\theta) = \beta_s\cos(\theta - 90^\circ) + \beta_c\cos(\theta - \theta_{\rm conf})$: an S-cone-axis component ($\beta_s$) plus a family-specific confusion-axis rotation ($\beta_c$; $\theta_{\rm conf} = 16^\circ$ protan, $150^\circ$ deutan). Outer ring, displayed hue; inner ring, hue perceived under the fitted distortion; labels give the signed displacement.
  \textbf{(2)} Three candidate loss atoms quantify CVD--HC differences. A three-gate procedure (direction, boundary saturation, held-out test-loss over $N = 300$ HC resamples; \S\ref{sec:methods:selection}) selected a different combination for each participant (deutan, $\gamma_{\rm OY} + L_{\rm RDM}^{(V2)}$; protan, $\gamma_{\rm all} + L_{\rm RDM}^{(V1)}$); the LOCO atom entered neither selected combination.
  \textbf{(3)} A grid search over $(\beta_s, \beta_c) \in [-90^\circ, 90^\circ]^2$ at $2^\circ$ resolution minimizes the $z$-scored composite to the selected fit $\hat\beta$. The surface shown is the deutan participant's; both landscapes and their resample uncertainty are in Figure~\ref{fig:landscape}.
  \textbf{(4)} The fitted forward map is numerically inverted: for each displayed hue $\theta$, the corrected stimulus $\theta_{\rm pre}$ solves $\theta_{\rm pre} + \delta\theta(\theta_{\rm pre}) = \theta$.
  \textbf{(5)} The per-subject filter is $\delta\theta_{\rm filter}(\theta) = \theta_{\rm pre} - \theta$; three of the eight hues are shown for the deutan participant, and the full eight-hue filter for both participants is in Figure~\ref{fig:filter}.
  Hue wheels and swatches are computed from each participant's fitted parameters; the schematic reports no new statistics.}
  \label{fig:pipeline}
\end{figure*}

\paragraph{Candidate Model 1: retinal-plus-gain (R+C).}
The prevailing account attributes CVD color distortion to a cone-spectral shift at the retinal level \cite{machado2009}, optionally augmented by a cortical compensation gain $g$ \cite{boehm2014,tregillus2021}:
%
\begin{equation}
  \delta\theta_{\rm RC}(\theta;\,g)
    = (2 - g)\,\delta\theta_{\rm Mach}(\theta)
  \label{eq:rc}
\end{equation}
%
where $\delta\theta_{\rm Mach}$ is the per-hue angular shift predicted by Machado simulation applied to Stockman \& Sharpe LMS cone fundamentals \cite{stockman2000}. The gain $g$ parameterizes cortical compensation: at $g = 1$ the cortical stage leaves the retinal distortion unchanged; at $g = 2$ it fully cancels the retinal shift; at $g > 2$ cortical processing reverses the direction of the retinal shift, moving the perceived hue past the undistorted color. The confusion axis is the color-space direction along which CVD observers cannot distinguish hues from one another. Because $\delta\theta_{\rm Mach}$ lies entirely along this axis, R+C can only displace colors along this single direction and cannot account for distortions in other color directions.

\paragraph{Candidate Model 2: two-component cortical (2-comp).}
We propose a parametric model of the hue distortion pattern measured in \S\ref{sec:results:loco} and \S\ref{sec:results:geometry}:
%
\begin{equation}
  \delta\theta(\theta;\,\beta_s,\beta_c)
    = \beta_s\cos(\theta - 90^\circ)
    + \beta_c\cos(\theta - \theta_{\rm conf})
  \label{eq:2comp}
\end{equation}
%
where $\theta_{\rm conf}$ is the CVD-subtype confusion axis ($16^\circ$ protan; $150^\circ$ deutan), derived from Stockman \& Sharpe cone fundamentals \cite{stockman2000}. $\beta_s$ parameterizes distortion along the S-cone axis ($\theta = 90^\circ$) and $\beta_c$ along the individual confusion axis. These two components parameterize the two color-space directions along which the distortion is hypothesized to lie. Both signs of $\beta_c$ are unconstrained. We constrain $\beta_s \geq 0$ as a weak directional prior that fixes only the sign, not the magnitude, of the S-cone-axis term, consistent with reports that CVD perceptual errors concentrate near the S-cone confusion loci \citeNP{emery2021}. Model adequacy was assessed by the procedure in \S\ref{sec:methods:selection}; results are reported in \S\ref{sec:results:twocomp}.

% ──────────────────────────────────────────────────────────────────────────────
\subsection{Inverse fitting}
\label{sec:methods:twocomp}
% ──────────────────────────────────────────────────────────────────────────────

Per-subject distortion parameters $(\beta_s,\beta_c)$ were estimated by finding the distortion that, when applied to HC-derived neural representations, best reproduces the observed CVD responses. Three loss families served as the fitting objective. A specific instance of a family --- for example, the aggregate JND ratio or the $\Delta$RDM loss at a single ROI --- is a \emph{loss atom}; candidate objectives are combinations of atoms, selected as described in \S\ref{sec:methods:selection}.

\paragraph{Behavioral JND ratio ($L_\gamma$).}
For each hue pair, the JND ratio $\gamma = \gamma_{\rm CVD} / \bar{\gamma}_{\rm HC}$ quantifies how much more difficult hue discrimination is for a CVD observer relative to the HC mean (\S\ref{sec:methods:behavioral}). Two variants enter the fitting: a subtype-targeted pair ratio $\gamma_{\rm subtype}$ (the hue pair on each participant's confusion axis predicted to show the largest distortion) and an aggregate ratio $\gamma_{\rm all}$ averaged across all tested pairs. $L_\gamma$ penalizes rank discordance between the model-predicted distortion magnitude at each pair (from Eq.\,\ref{eq:2comp}) and the observed $\gamma$ ordering.

\paragraph{Representational $\Delta$RDM cosine ($L_{\rm RDM}$).}
For a candidate distortion, a simulated difference $\Delta\text{RDM}_{\rm sim}$ is obtained by applying $\delta\theta$ to the HC mean RDM and comparing it to the observed $\Delta\text{RDM}$ (\S\ref{sec:methods:rdm}). The loss is the cosine dissimilarity between the observed and simulated 28-element upper-triangle vectors:
%
\begin{equation}
  L_{\rm RDM} = \frac{1 - \cos\!\bigl(\Delta\text{RDM}_{\rm sim},\,
    \Delta\text{RDM}_{\rm obs}\bigr)}{2}
  \label{eq:lrdm}
\end{equation}

\paragraph{hV4 LOCO voxel-prediction ($L_{\rm LOCO}$).}
The per-hue vulnerability profile $\mathbf{v}_{\rm obs} \in [0,1]^8$ records where each CVD participant's hue interpolation breaks down (\S\ref{sec:methods:loco}). For a candidate $(\beta_s,\beta_c)$, the simulated profile $\mathbf{v}_{\rm sim}$ is obtained by re-applying the HC-trained ridge-GCV encoding weights $W$ (\S\ref{sec:methods:encoding}) to the channel design matrix $C$ shifted by the candidate distortion. The loss is the mean squared deviation between the two profiles:
%
\begin{equation}
  L_{\rm LOCO} = \frac{1}{8}\,\|\mathbf{v}_{\rm sim} -
    \mathbf{v}_{\rm obs}\|^2
  \label{eq:lloco}
\end{equation}

\paragraph{Grid search.}
Parameters were searched by exhaustive grid at $2^\circ$ resolution, with $\beta_s$ over $[0^\circ, 50^\circ]$ (26 values) and $\beta_c$ over $[-50^\circ, 50^\circ]$ (51 values), giving $1{,}326$ cells. For R+C (Eq.\,\ref{eq:rc}), $g$ was searched over $[0, 3]$ in steps of $0.05$ (61 values). Per-subject loss combinations and parameter selection are described in \S\ref{sec:methods:selection}.

% ──────────────────────────────────────────────────────────────────────────────
\subsection{Parameter selection}
\label{sec:methods:selection}
% ──────────────────────────────────────────────────────────────────────────────

Per-subject loss combinations and model class were selected by a three-gate procedure applied to results from a 5-train/2-test HC resample ($N = 300$) and a strict 7-fold leave-one-HC-out using the same loss atoms.

\paragraph{Gate 1: Directional precondition.}
Only loss atoms where the CVD loss value exceeded the HC LOO distribution in the expected direction were admitted. Admission required signed Cohen's $d \geq +0.5$; atoms failing this gate were excluded from combination.

\paragraph{Gate 2: Boundary saturation.}
Among admitted combinations, those where $\geq 50\%$ of resample
solutions saturated the grid boundary were rejected as indicative of model misspecification.

\paragraph{Gate 3: Held-out test-loss.}
The primary criterion was the median composite test-loss ($\overline{L}_{\rm test}$) across the $N = 300$ HC resamples. For each resample, the composite loss was evaluated at the training argmin on the two held-out HC, z-normalized relative to the training grid. Lower $\overline{L}_{\rm test}$ indicates generalization to unseen HC references. The secondary criterion was the test-loss IQR across resamples, measuring resample stability of the generalization estimate. A dedicated neural RDM held-out loss ($\overline{L}_{\rm RDM}^{\rm LOO}$) from the strict 7-fold leave-one-out isolates the neural component's held-out performance from the behavioral weighting in the composite. Supplementary stability indices included the IQR of recovered $(\hat\beta_s, \hat\beta_c)$ (parameter stability), the fraction of resamples recovering the modal $45^\circ$ bin (mode share), and the strict 7-fold LOO parameter range. The HC false-positive rate was not itself used as a selection criterion; it is reported descriptively in \S\ref{sec:results:identifiability}.

% ──────────────────────────────────────────────────────────────────────────────
\subsection{Identifiability and recovery}
\label{sec:methods:identifiability}
% ──────────────────────────────────────────────────────────────────────────────

Four pre-specified checks were run on each production candidate before any parameter was reported. (1)~Parameter recovery: synthetic CVD data at a known ground truth were forward-projected through each HC donor's encoder with spatially correlated voxel noise, and the full fitting pipeline was re-run ($n = 140$ per candidate); recovery was summarized by $f_{10^\circ}$, the fraction of refits within $10^\circ$ of ground truth on both axes (pass $f_{10^\circ} \geq 0.5$, $|\text{bias}| < 10^\circ$). (2)~Origin-recovery null: the same synthesis with ground truth at $(0^\circ, 0^\circ)$ and each donor's real JND, recording the procedure's per-axis uncertainty floor (excluded from the test pool, as it returns no $p$-value). (3)~HC pseudo-CVD specificity: each HC participant's real amplitudes were substituted into the CVD slot in turn, giving a seven-member null and a one-sided rank-distance percentile. (4)~Color-label permutation: trial labels of the real CVD data were shuffled $N = 1000$ times with the HC reference held fixed, giving a one-sided $p_{\rm perm}$. The three test-bearing checks across the two production candidates (six tests) were corrected together under the Benjamini--Hochberg false-discovery rate ($\alpha = 0.05$); outcomes are reported in \S\ref{sec:results:identifiability}.

% ──────────────────────────────────────────────────────────────────────────────
\subsection{Stimulus-space filter}
\label{sec:methods:filter}
% ──────────────────────────────────────────────────────────────────────────────

The subject-specific correction filter was derived as the pre-image of the fitted 2-component transform $T(\hat{\beta}_s, \hat{\beta}_c): \theta \mapsto \theta + \delta\theta(\theta;\hat{\beta}_s,\hat{\beta}_c)$. For each of the 8 displayed hues $\theta_k$, we solved for the input angle $\tilde{\theta}_k$ such that $T(\tilde{\theta}_k) = \theta_k$, yielding a per-hue correction $\delta\theta_k^{\rm filt} = \tilde{\theta}_k - \theta_k$. A CVD observer viewing $\tilde{\theta}_k$ thus receives, under the model, the same cortical hue representation as a healthy control observer viewing $\theta_k$.

Pre-images were computed by numerical inversion of the forward map (Brent's-method root-finding) to a convergence tolerance of $< 0.001^\circ$.

% ──────────────────────────────────────────────────────────────────────────────
\subsection{Filter evaluation}
\label{sec:methods:filter_eval}
% ──────────────────────────────────────────────────────────────────────────────

Both CVD participants completed a second scanning session that compared the individualized filter and a deployed macOS accessibility filter against the unfiltered Session-1 baseline. The individualized filter was the per-subject 2-component stimulus-space pre-image, with parameters frozen from the main analysis and therefore out-of-sample. Acquisition, comparator, run-count, and single-case-inference details are given in Supplementary~\S\ref{app:filter_eval}.

On the second-session data we recomputed, within each condition, the primary Session-1 interpolation metric (LOCO adjacent accuracy, \S\ref{sec:methods:loco}), LORO classification, SRM Procrustes disparity into the HC shared space, and the representational dissimilarity matrices. As a secondary encoding index we also computed the forward-tuning correlation, the Spearman $\rho$ between each ROI's per-hue decoded hue and the canonical hue map. Because it indexes the encoding fit rather than the validated decoding readout, it is reported as corroboration only.

Each condition's neural index was reported as a standardized single-case effect size (Cohen's $d$ against the HC distribution) rather than an inferential $p$, because the grand-mean permutation null is biased for this design. In both CVD participants, the JND and 8AFC tasks (\S\ref{sec:methods:behavioral}) were repeated under each filter, each compared to the unfiltered Session-1 baseline and to the HC distribution; this interleaved, cross-session behavioral comparison is descriptive.

% ──────────────────────────────────────────────────────────────────────────────
\subsection{Reproducibility}
\label{sec:methods:repro}
% ──────────────────────────────────────────────────────────────────────────────

All analyses were conducted in Python 3.10 (numpy 1.24.3, scipy 1.11.3, scikit-learn 1.3.0, BrainIAK 0.11). Random seeds were fixed (\texttt{seed=42}) for all permutation tests and bootstrap resampling. Analysis code and preprocessing scripts are available at \url{https://github.com/Transconnectome/colorBlind_analysis}. Anonymized preprocessed data are available on reasonable request to the corresponding author.
